\documentclass[a4paper, 11pt]{article}
\usepackage[brazilian]{babel}
\usepackage[T1]{fontenc}
\usepackage[utf8]{inputenc}
\usepackage{graphicx}
\usepackage{float}
\usepackage{hyperref}
\usepackage{tikz}
\usepackage{pgfgantt}

\usepackage{setspace}

\usepackage{listings}
\usepackage{xcolor}

\graphicspath{ {./tables} }

\definecolor{pinkglamour}{HTML}{ff7675}
\definecolor{greendarnertail}{HTML}{74b9ff}
\definecolor{shymoment}{HTML}{a29bfe}
\definecolor{lightgreenishblue}{HTML}{55efc4}
\definecolor{picopink}{HTML}{fd79a8}

\usepackage{booktabs}
\usepackage{fontspec}
\usepackage{xeCJK}
\setCJKmainfont{Noto Serif CJK JP}

\makeatletter

\renewcommand{\lstlistingname}{Código}
\renewcommand{\lstlistlistingname}{Lista de \lstlistingname s}

\lstdefinestyle{mystyle}{
    language = [Sharp]C,
    % Colors
    keywordstyle    = \color{violet},
    stringstyle     = \color{red},
    commentstyle    = \color{gray}\ttfamily\footnotesize,
    backgroundcolor = \color{white},
    rulesepcolor    = \color{darkgray},
    identifierstyle = \color{black},
    % Caracteres
    basicstyle       = \ttfamily\footnotesize,
    extendedchars    = true,
    showspaces       = false,
    showstringspaces = false,
    % Frame Style
    numbers         = left,
    tabsize         = 2,
    numberstyle     = \scriptsize,
    rulecolor=\color{black},
    frame           = single,
    breaklines      = true,
    breakautoindent = true,
    framexleftmargin     = 2.5em,
    aboveskip       = 10pt,
    belowskip       = 10pt,
    belowcaptionskip = 5pt,
    captionpos      = b,
    % Caracteres especiais
    literate =
    {á}{{\'a}}1 {é}{{\'e}}1 {í}{{\'i}}1 {ó}{{\'o}}1 {ú}{{\'u}}1
    {Á}{{\'A}}1 {É}{{\'E}}1 {Í}{{\'I}}1 {Ó}{{\'O}}1 {Ú}{{\'U}}1
    {à}{{\`a}}1 {è}{{\`e}}1 {ì}{{\`i}}1 {ò}{{\`o}}1 {ù}{{\`u}}1
    {À}{{\`A}}1 {È}{{\`E}}1 {Ì}{{\`I}}1 {Ò}{{\`O}}1 {Ù}{{\`U}}1
    {Ä}{{\"A}}1 {Ë}{{\"E}}1 {Ï}{{\"I}}1 {Ö}{{\"O}}1 {Ü}{{\"U}}1
    {â}{{\^a}}1 {ê}{{\^e}}1 {î}{{\^i}}1 {ô}{{\^o}}1 {û}{{\^u}}1
    {Â}{{\^A}}1 {Ê}{{\^E}}1 {Î}{{\^I}}1 {Ô}{{\^O}}1 {Û}{{\^U}}1
    {ã}{{\~a}}1 {ẽ}{{\~e}}1 {ĩ}{{\~i}}1 {õ}{{\~o}}1 {ũ}{{\~u}}1
    {Ã}{{\~A}}1 {Ẽ}{{\~E}}1 {Ĩ}{{\~I}}1 {Õ}{{\~O}}1 {Ũ}{{\~U}}1
    {ő}{{\H{o}}}1 {Ő}{{\H{O}}}1
    {ç}{{\c c}}1 {Ç}{{\c C}}1
}

\lstset{style=mystyle}

\def\ps@myheadings{
    \let\@oddhead\@empty
    \let\@evenhead\@empty
    \def\@oddfoot{\hspace*{-2cm}\hfill\thepage\hspace*{-2cm}}%
    \def\@evenfoot{\hspace*{-2cm}\thepage\hfill\hspace*{-2cm}}%
}

\newcommand{\subtitle}[1]{\gdef\@subtitle{#1}}


\def\@maketitle{%
    \newpage
    \null
    \vskip -50em%
    \begin{center}%
    {\Large
            Instituto de Matemática, Estatística e Computação\\%
                        Universidade de São Paulo\\%
                        \par
    }
    \vskip 2.5em%
    \let \footnote \thanks
    {\LARGE \@subtitle \par}%
    \vskip 2.5em%
    {\LARGE \textbf{\@title} \par}%
    \vskip 1.5em%
    {\large
      \lineskip .5em%
      \begin{tabular}[t]{c}%
        \@author
      \end{tabular}\par}%
    \vskip 10em%
    \end{center}%
    \par
    \vskip 1.5em
    }

\title{Desenvolvimento de Ferramenta de Ensino de Shogi com Feedback Adaptativo para Iniciantes}
\subtitle{Trabalho de Conclusão de Curso}

\author{Autor: Vinícius Yuiti Okuma Shimizu\\
        Orientador: Denis Deratani Mauá}
\date{}

\hypersetup{
    colorlinks = true,
    allcolors = black,
    urlcolor = blue,
    pdftitle = {Desenvolvimento de ferramenta de ensino adaptativo de shogi com feedback para iniciantes - Trabalho de Conclusão de Curso},
    pdfauthor = {Vinícius Yuiti Okuma Shimizu},
    pdfcreator = {\LaTeX},
}

\def\Hline{%
\noalign {\ifnum0 =`}\fi \hrule \@height 1.5pt \futurelet \reserved@a \@xhline}

\makeatother

\pagestyle{myheadings}
\begin{document}
\onehalfspace
    \maketitle\thispagestyle{myheadings}
    \newpage

    \tableofcontents

    % \newpage
    % \section{Resumo}
    % \paragraph{}
    % Este trabalho teve como objetivo desenvolver uma plataforma de ensino de Shogi voltada para iniciantes, utilizando como metodologia o aprendizado por reforço aliado a exercícios
    % sugeridos com base no desempenho do usuário e feedback gerado por um modelo de LLM (\textit{Large Language Model}). Os exercícios foram gerados automaticamente a partir de partidas reais e
    % cobrem os fundamentos do jogo, desde a identificação das peças até o chequemate.
    \newpage
    \section{Introdução}
    \paragraph{}
    Shogi é um jogo de estratégia japonês semelhante ao xadrez, mas que apresenta 
    mecânicas adicionais que aumentam a dificuldade de seu aprendizado. Entre essas mecânicas,
    destacam-se a possibilidade de reutilizar peças capturadas e a promoção de peças para versões com
    movimentos diferentes. 
    \paragraph{}
    Ao longo dos anos, plataformas como o Lishogi surgiram com o intuito 
    de facilitar o acesso ao aprendizado, oferecendo recursos como partidas online ou resolução de problemas como 
    formas de treino. Essas plataformas, no entanto, apresentam limitações que se tornam uma barreira ao aprendizado 
    de iniciantes. O Lishogi, por exemplo, disponibiliza exercícios simples para quem está aprendendo, 
    mas não fornece informações sobre os motivos pelos quais uma resposta está incorreta. O mesmo vale para 
    as análises de partidas, nas quais as jogadas recebem uma pontuação de acordo com sua qualidade, sem explicar por que
    determinada decisão seria melhor.
    \paragraph{}
    Para um iniciante que deseja aprender a jogar de forma autônoma, a falta de orientação pode dificultar
    a compreensão dos conceitos envolvidos e assim limitar o aproveitamento dos exercícios realizados. 
    Nesse contexto, o uso de modelos de linguagem pode contribuir para suprir tal necessidade ao fornecer ao jogador explicações
    moldadas ao seu nível de conhecimento, desempenhando um papel semelhante ao de um tutor. 
    \paragraph{}
    Diante desse cenário, este trabalho propõe o desenvolvimento do Shogi Sensei, uma plataforma web voltada ao ensino
    introdutório de shogi. O projeto fornece exercícios divididos em módulos que abordam os conceitos fundamentais
    do jogo: reconhecimento das peças, movimentação, promoção, reintrodução de peças capturadas e situações de 
    xeque-mate.
    Ao final de cada lista, o usuário recebe os resultados, incluindo as questões respondidas incorretamente e uma análise
    gerada por um modelo de linguagem (LLM) que explica os motivos do erro. A seleção de exercícios de cada lista é adaptada
    ao desempenho do jogador, aumentando a proporção de exercícios dos módulos nos quais ele apresenta maior dificuldade.
    Para avançar para o próximo módulo, é necessário atingir uma média mínima nas listas mais recentes do módulo atual.
    \paragraph{}
    A avaliação do sistema foi realizada durante 2 semanas com X usuários. Durante esse período, 
    foram registradas métricas como taxa de acerto, tempo de resposta e distribuição de exercícios por módulo para
    acompanhar o desempenho dos usuários ao longo do experimento.
    \paragraph{}
    O restante deste trabalho está organizado da seguinte forma. O capítulo 2 apresenta os conceitos e tecnologias
    utilizados para desenvolver este trabalho. O capítulo 3 apresenta trabalhos relacionados. O capítulo 4 descreve a implementação
    do sistema. O capítulo 5 contém a metodologia utilizada para avaliar a plataforma desenvolvida. O capítulo 6 
    apresenta e discute os resultados obtidos no experimento descrito no capítulo anterior. Por fim, o capítulo 7
    apresenta as conclusões e possíveis implementações futuras. 
    \section{Conceitos e tecnologias}
    \subsection[Shogi]{Shogi}
    O shogi é um jogo de tabuleiro disputado em um tabuleiro com 9 linhas e 9 colunas. O jogo utiliza 8 tipos de peças, das quais 6 podem assumir formas promovidas (que possuem movimentações diferentes das versões padrão), totalizando 14 tipos de peças. A notação ocidental das peças e seus respectivos símbolos japoneses são apresentados na Tabela~\ref{tab:shogi_pieces}.
    \paragraph{}
    A notação ocidental utilizada em publicações de língua inglesa foi desenvolvida por George Hodges e Glyndon Townhill em 1976 \textbf{[COLOCAR REFERENCIA]}, com base em elementos da notação algébrica do xadrez. 
    \paragraph{Convenções:}
    \begin{table}[H]
        \centering
        \caption{Correspondência entre as notações japonesa e ocidental das peças de shogi}
        \label{tab:shogi_pieces}
        \begin{tabular}{c c c}
            \toprule
            \textbf{Símbolo} & \textbf{Notação ocidental} & \textbf{Peça} \\
            \midrule
            歩 & P & Peão \\
            香 & L  & Lança \\
            桂 & N  & Cavalo \\
            銀 & S  & General de Prata \\
            金 & G  & General de Ouro \\
            角 & B  & Bispo \\
            飛 & R  & Torre \\
            玉 & K  & Rei \\
            と & +P  & Peão Promovido \\
            成香 & +L  & Lança Promovida \\
            成桂 & +N  & Cavalo Promovido \\
            成銀 & +S  & General de Prata Promovido \\
            馬 & +B  & Bispo Promovido \\
            龍 ou 竜 & +R  & Torre Promovida \\
            \bottomrule
        \end{tabular}
    \end{table}
    \begin{itemize}
        \item \textbf{Jogadores}: os jogadores são definidos como preto (sente), quem faz o primeiro movimento, ou branco (gote).
        \item \textbf{Casas}: as colunas são numeradas de 1 a 9, da direita para a esquerda, e as linhas são identificadas com letras de a a i, de cima para baixo, considerando o ponto de vista do sente.
        \item \textbf{Posição}: uma casa é identificada pela combinação do número da coluna com a letra da linha. Por exemplo, 5a corresponde à casa localizada na coluna 5 e na primeira linha.
        \item \textbf{Movimentos}\textbf{[INCLUIR REFERENCIA PARA REGRAS]}: o sistema adota uma variação da convenção de George, com adaptações para facilitar a compreensão por iniciantes. Os movimentos de cada peça se encontram na Tabela~\ref{tab:pieces_movements}.
            \begin{itemize}
                \item \textbf{Deslocamento}: [Peça][Origem]->[Destino]. Por exemplo, \textbf{P9g->9f} se refere ao movimento de um peão da casa \textbf{9g} para a casa \textbf{9f}. Na notação de referência, o deslocamento não possui o símbolo "->".
                \item \textbf{Captura}: quando aplicável, o movimento ocupa a casa da peça adversária capturada, que é removida do tabuleiro e passa a pertencer à mão jogador que realizou a captura.
                \item \textit{\textbf{Drop}}: [Peça]*[Posição]. Refere-se à recolocação no tabuleiro de uma peça capturada anteriormente. Por exemplo, \textbf{P*9e} representa a colocação de um peão na casa \textbf{9e}.
                \item \textbf{Promoção}: [Peça][Origem]->[Destino]+. O símbolo "+" representa que a peça foi promovida ao final do movimento. Por exemplo, \textbf{P9d->9c+} representa o movimento do peão da casa \textbf{9d} para a casa \textbf{9c}, seguido de sua promoção. A promoção pode ocorrer quando o movimento começa, termina ou acontece inteiramente dentro da zona de promoção, que corresponde às três fileiras mais distantes do jogador. Em determinadas situações ela é obrigatória: por exemplo, um peão ou uma lança que alcança a última fileira, ou um cavalo que alcança uma das duas últimas fileiras, pois não possuem mais movimentos possíveis a partir dessas posições. 
            \end{itemize}
    \end{itemize}
    \begin{table}[H]
        \centering
        \caption{Movimentos das peças}
        \label{tab:pieces_movements}
        \begin{tabular}{c p{9cm}}
            \toprule
            \textbf{Peça} & \textbf{Movimento}\\
            \midrule
            Peão & Uma casa para a frente. \\
            Lança & Qualquer número de casas para a frente. \\
            Cavalo & Salta para uma das duas casas duas fileiras à frente e uma coluna para um dos lados. \\
            General de prata & Uma casa na diagonal ou uma casa para a frente. \\
            General de ouro & Uma casa para a frente, para trás ou para os lados, ou uma casa na diagonal para a frente. \\
            Bispo & Qualquer número de casas na diagonal. \\
            Torre & Qualquer número de casas na horizontal ou vertical. \\
            Rei & Uma casa em qualquer direção. \\
            Peão promovido & Como um general de ouro. \\
            Lança promovida & Como um general de ouro. \\
            Cavalo promovido & Como um general de ouro. \\
            General de prata promovido & Como um general de ouro. \\
            Bispo promovido & Como um bispo ou uma casa na horizontal ou vertical. \\
            Torre promovida & Como uma torre ou uma casa na diagonal. \\
            \bottomrule
        \end{tabular}
    \end{table}
    \section{Trabalhos relacionados}
    \section{Desenvolvimento}
    \section{Metodologia}
    \section{Resultados}
    \section{Conclusão}


    \cleardoublepage\thispagestyle{empty}
    \bibliographystyle{abnt}
    \bibliography{references}
\end{document}
